\section{Exact majority response}
\label{hrec:app:majority}
With \(p=(1+v)/2\), the mean-majority polynomial is
\begin{equation}
 B_N(v)=\sum_{j=0}^N
 \operatorname{sign}(2j-N)\binom Nj p^j(1-p)^{N-j}.
\end{equation}
At \(u=0\), each spin configuration has probability \(2^{-N}\). Differentiation using Eq.~\eqref{hrec:eq:detector} multiplies that probability by \(\beta S_N\), giving the first line of Eq.~\eqref{hrec:eq:exact-slope}. For the binomial sum, use
\[
 (2j-N)\binom Nj
 =N\left[\binom{N-1}{j-1}-\binom{N-1}j\right].
\]
Summing over \(j>N/2\) telescopes. Doubling the positive contribution by symmetry yields
\(\E|S_N|=N2^{-(N-1)}\binom{N-1}{\lfloor(N-1)/2\rfloor}\).
A tie contributes zero to this derivative.

For \(N=2j\), Stirling expansion of the central binomial coefficient gives
\(\E|S_N|=\sqrt{2N/\pi}[1-1/(4N)+O(N^{-2})]\).
For \(N=2j+1\), the corresponding expansion has \(+1/(4N)\).
For even \(N\), the literal-zero tie rule has variance
\(1-\Pr(S_N=0)\) at zero bias, while a fair binary tie-break has variance one. Their equal mean laws therefore do not license an identical binary stochastic recursion.

